% ---------------------------------------------------------------------------------------
% MODELO DE TCC - INSTITUTO FEDERAL DE SERGIPE (IFS)
% ---------------------------------------------------------------------------------------
% Autor do Modelo: Mikael dos Santos Gois (mikaelgois1@gmail.com)
% Baseado nos modelos de: 
%   - Diogo D. Moreira (TCC, Dissertação e Tese)
%   - DECOM/UFOP (Monografia de Ciência da Computação)
% Modelo Base: abnTeX2 adaptado para as normativas internas de TCC do IFS (Art. 31 e 34)
% Versão: 1.0 (2026)
% ---------------------------------------------------------------------------------------
% Este arquivo faz parte de um modelo opensource para estudantes do IFS.
% Não remova os créditos dos autores originais e do modelo base.
% Este modelo é fornecido "no estado em que se encontra", sem garantias de qualquer tipo.
% ---------------------------------------------------------------------------------------

% Modelo de ficha catalográfica. Você deverá substituir esta
% folha na versão final da monografia por um pdf fornecido pela biblioteca.
% Salve o modelo oficial como ficha_catalografica.pdf na raiz do projeto,
% e use o comando abaixo para inseri-lo na versão final do texto.

%\begin{fichacatalografica}
%    \includepdf{ficha_catalografica.pdf}
%\end{fichacatalografica}

%% Modelo de Como fazer a Folha/Ficha Catalográfica:

\begin{fichacatalografica}
  \sffamily
  \vspace*{\fill}					% Posição vertical
  \begin{center}					% Minipage Centralizado
    \fbox{
      \begin{minipage}[c][8cm]{14cm}		% Largura
        \small
        \begin{center}
          {Isto é apenas um exemplo e deve ser substituído pela ficha catalográfica oficial fornecida pela biblioteca do IFS.}
        \end{center}
        \imprimirautorcite.
        %Sobrenome, Nome do autor

        \hspace{0.5cm}  \\
        \imprimirtitulo. / \imprimirautor. -- \imprimirlocal, \imprimirano.

        % \hspace{0.5cm} \pageref{LastPage} p. 1 :il. (colors; grafs; tabs).\\

        \hspace{0.5cm}  \\
        {Trabalho de conclusão de curso (Bacharelado em Ciência da
        Computação) – Instituto Federal de Educação, Ciência e Tecnologia de
        Sergipe – IFS, \imprimirano. \\}

        \hspace{0.5cm} \imprimirorientadorRotulo~\imprimirorientador.
        \ifthenelse{\equal{\imprimircoorientador}{}}{}{ \imprimircoorientadorRotulo~\imprimircoorientador \\}

        % \hspace{0.5cm}
        % \parbox[t]{\textwidth}{\imprimirtipotrabalho~--~\imprimirinstituicao, ~ \\
        %   \imprimirinstituto, ~\imprimirdepartamento,~\imprimirano.}\\

        \hspace{0.5cm} % Palavras-chave do trabalho
        1. Palavra-chave.
        2. Palavra-chave.
        2. Palavra-chave.
        3. Palavra-chave.
        4. Palavra-chave.
        5. Palavra-chave.
        I. IFS.
        II. \imprimirautorcite.
        III. \imprimirorientadorcite.
        \ifthenelse{\equal{\imprimircoorientadorcite}{}}{IV. \imprimirtitulo.}{IV. \imprimircoorientadorcite. V. \imprimirtitulo.}
      \end{minipage}
    }
  \end{center}
\end{fichacatalografica}
\addtocounter{page}{-1} % Subtrai 1 da contagem total de páginas
